\documentclass[11pt]{article}
\usepackage[T1]{fontenc}
\usepackage[utf8]{inputenc}
\usepackage{lmodern}
\usepackage[ngerman]{babel}
\usepackage{amsmath,amssymb,mathtools,array,booktabs,pdflscape,adjustbox,diagbox}
\usepackage{geometry}
\usepackage{microtype}
\geometry{a4paper,margin=1.45cm}
\setlength{\parindent}{1.25em}
\setlength{\parskip}{0pt}
\makeatletter
\renewcommand\footnotesize{\@setfontsize\footnotesize{7.5}{8.5}\selectfont}
\makeatother
\allowdisplaybreaks
\setlength{\emergencystretch}{10em}
\sloppy
\hbadness=10000
\vbadness=10000
\hfuzz=2pt
\vfuzz=10pt
\raggedbottom

% current section packet macros
\providecommand{\Sp}{\operatorname{Sp}}
\providecommand{\Tr}{\operatorname{Tr}}
\providecommand{\Norm}{\operatorname{N}}
\providecommand{\Mat}{\operatorname{Mat}}
\providecommand{\Gal}{\operatorname{Gal}}
\providecommand{\Cl}{\operatorname{Cl}}
\providecommand{\Gg}{\mathfrak G}
\providecommand{\Hh}{\mathfrak H}
\providecommand{\Ii}{\mathfrak I}
\providecommand{\Jj}{\mathfrak J}
\providecommand{\Aa}{\mathfrak A}
\providecommand{\Bb}{\mathfrak B}
\providecommand{\Cc}{\mathfrak C}
\providecommand{\Oo}{\mathfrak O}
\providecommand{\oo}{\mathfrak o}
\providecommand{\pp}{\mathfrak p}
\providecommand{\PP}{\mathfrak P}
\providecommand{\LL}{\mathfrak L}
\providecommand{\RR}{\mathfrak R}
\providecommand{\ZZ}{\mathfrak Z}
\providecommand{\ee}{\mathfrak e}

\begin{document}

\clearpage

\section*{41. Der Hauptgeschlechtssatz für relativ-galoissche Zahlkörper}

\begin{center}
\emph{Math. Ann. 108 (1933), S. 411--419}
\end{center}

In neuester Zeit hat es sich gezeigt, daß die nichtkommutativen Methoden, insbesondere die Theorie der Algebren, die Möglichkeit geben, bekannte Sätze über relativ-zyklische und relativ-abelsche Zahlkörper für beliebige relativ-galoissche Körper zu formulieren und zu beweisen.\footnote{Vgl. dazu meinen Züricher Vortrag, \emph{Hyperkomplexe Systeme in ihren Beziehungen zur kommutativen Algebra und Zahlentheorie}, Verhandlungen des Internationalen Mathematiker-Kongresses Zürich 1932, Bd. I. Für den Satz über zerfallende Algebren vgl. außerdem H. Hasse, \emph{Die Struktur der R. Brauerschen Algebrenklassengruppe}, Math. Ann. 107 (1932).} Ich erinnere vor allem an den Normensatz, der sich allgemein als ein Satz über zerfallende Algebren ausspricht. Im folgenden zeige ich, daß sich entsprechend auch der Hauptgeschlechtssatz allgemein hyperkomplex formulieren läßt, als ein Satz über innere Automorphismen oder auch über verschränkte Darstellungen. Der Beweis ergibt sich aus dem genannten Satz über zerfallende Algebren durch rein algebraisch-arithmetische Überlegungen, während im ersteren Satz der Normensatz im zyklischen Fall, und schon bei Primzahlgrad, als transzendenter Kern steckt.

Zum besseren Verständnis der Formulierung schicke ich den im folgenden nicht gebrauchten Hauptgeschlechtssatz im Minimalen voraus,\footnote{Ich sage ``im Minimalen'', da es sich um das algebraische Analogon des Hauptgeschlechtssatzes in beliebigen Körpern handelt, während im ``Kleinen'' schon die Bedeutung des Übergangs zum \(p\)-adischen Grundkörper hat.} einen elementar-algebraischen Satz,\footnote{Der Satz findet sich bei A. Speiser, Zahlentheoretische Sätze aus der Gruppentheorie, Math. Z. 5 (1919), S. 1--6; er ist dort bereits als Satz über verschränkte Darstellungen ausgesprochen.} der im zyklischen Spezialfall in den bekannten Satz 90 aus Hilberts Zahlbericht übergeht, wonach \(\Norm(a)\) dann und nur dann gleich Eins wird, wenn \(a=b^{1-S}\). Und zwar spreche ich auch schon diesen Satz als Satz über innere Automorphismen oder über verschränkte Darstellungen aus und beweise ihn vermöge der allgemeinen Sätze über innere Automorphismen und verschränkte Produktdarstellung der Algebren.

An Stelle der letzteren tritt beim eigentlichen Hauptgeschlechtssatz das verschränkte Produkt von Idealklassengruppe und galoisscher Gruppe, aber mit einer feineren induzierten Idealklasseneinteilung für die Faktorensysteme, wodurch ein Analogon zur Strahlklasseneinteilung der Klassenkörpertheorie erfaßt wird. Da hier allgemeine Sätze über Automorphismen und Darstellungen noch fehlen, kann der Hauptgeschlechtssatz als ein erster Schritt in dieser Richtung aufgefaßt werden. Ich gebe deshalb, wie erwähnt, eine Zurückführung vermöge des Satzes über die zerfallenden Algebren. Dadurch wird der Satz auf den entsprechenden für Ideale statt Idealklassen zurückgeführt; hier können die Brandtschen Sätze über den Zusammenhang der Maximalordnungen einer Algebra\footnote{Die Brandtsche Theorie der Maximalordnungen findet sich dargestellt und neu begründet bei H. Hasse, Über \(p\)-adische Schiefkörper und ihre Bedeutung für die Arithmetik hyperkomplexer Zahlsysteme, Math. Ann. 104 (1931).} als Äquivalent der Automorphismensätze aufgefaßt werden. Sie geben im Zusammenhang mit Betrachtungen über Maximalordnungen verschränkter Produkte\footnote{Je eine Note von C. Chevalley, H. Hasse und E. Noether wurde angekündigt; die beiden letzteren in einem Herbrand-Gedächtnisband.} einen dem Beweis im Minimalen wesentlich parallel laufenden, wenn auch durch Zurückgehen auf die einzelnen Stellen komplizierteren Beweis. Ich ziehe daher den von E. Artin angegebenen fast trivialen Beweis vor, der ein hier noch einfacheres Analogon des Speiserschen Beweises darstellt.

\subsection*{§1. Hauptgeschlechtssatz im Minimalen}

In diesem Paragraphen bedeute \(k\) einen beliebigen kommutativen Körper, nicht notwendig einen Zahlkörper; \(K/k\) sei eine separable galoissche Erweiterung \(n\)-ten Grades, \(\Gg\) die zugehörige Galoissche Gruppe.

Zunächst seien die bekannten Tatsachen über verschränkte Produkte zusammengestellt.\footnote{Die Theorie ist im Anschluß an eine Vorlesung von mir (1929/30) dargestellt bei H. Hasse, \emph{Theory of cyclic algebras}, Trans. Amer. Math. Soc. 34 (1932), Kap. 2; vgl. ferner M. Deuring über hyperkomplexe Zahlen.} Das verschränkte Produkt bedeutet eine gleichzeitige Einbettung von \(K\) und \(\Gg\) in eine Algebra \(A\) derart, daß die Automorphismen von \(K\) innere werden. Sind \(u_{S_1},\ldots,u_{S_n}\) Symbole zu den \(n\) Gruppenelementen, so setzt man \(A\) zunächst als Linearformenmodul vom Rang \(n\) über \(K\) an:
\[
        A=u_{S_1}K+\cdots+u_{S_n}K .                         \tag{1}
\]
Vermöge der Forderung des inneren Automorphismus, erzeugt durch \(u_S\), allgemeiner durch \(u_S K^*\), wird \(A\) zu einem Ring, also zu einer Algebra vom Rang \(n^2\) über \(k\). Die Forderung drückt sich aus durch
\[
        u_S^{-1}zu_S=z^S,\qquad\text{oder}\qquad
        z u_S=u_S z^S,                                      \tag{2}
\]
für jedes \(z\in K\), sowie durch
\[
        u_Su_T=u_{ST}a_{S,T},\qquad a_{S,T}\in K^*,           \tag{3}
\]
und durch das Assoziativgesetz
\[
        a_{S,T}^{\,R}a_{ST,R}=a_{S,TR}a_{T,R}.                \tag{4}
\]
Definiert man das Produkt zweier beliebiger Elemente durch
\[
 \Bigl(\sum_S u_S b_S\Bigr)\Bigl(\sum_T u_T c_T\Bigr)
   =\sum_{S,T}u_{ST}\,a_{S,T}\,b_S^{\,T}c_T ,
\]
so wird \(A\) zum verschränkten Produkt von \(K\) mit \(\Gg\) bei Faktorensystem \(a_{S,T}\); Bezeichnung
\[
        A=(a_{S,T},K,\Gg).
\]
Man beweist, daß \(A\) eine einfache normale Algebra über \(k\) wird, also ein Matrizenring \(D_r\) \(r\)-ten Grades über der zugeordneten Divisionsalgebra \(D\), und daß \(K\) maximaler kommutativer Unterkörper, also Zerfällungskörper ist. Umgekehrt gibt es zu jeder solchen normalen Divisionsalgebra \(D\) Matrizenringe \(D_r\), die auf diese Weise als verschränkte Produkte erzeugbar sind.

Geht man von \(u_S\) zu \(v_S=u_Sc_S\) mit \(c_S\in K^*\) über, so entstehen assoziierte Faktorensysteme
\[
        \bar a_{S,T}
        =a_{S,T}\frac{c_S^{\,T}c_T}{c_{ST}} .                 \tag{5}
\]
Insbesondere werden nach (5) die Faktorensysteme
\[
        \frac{c_S^{\,T}c_T}{c_{ST}}
\]
zu Eins assoziiert; diese Größen heißen Transformationsgrößen. Assoziierte Faktorensysteme werden in eine Klasse \((a)\) zusammengefaßt. Die Klassen \((a)\) entsprechen eineindeutig den direkten Produktbildungen einer abelschen Gruppe; sie bilden gegenüber direkter Produktbildung eine abelsche Gruppe, isomorph dem gliedweisen Produkt der Klassen von Faktorensystemen. Das Einselement ist die zerfallende Algebrenklasse, d. h. das System aller Transformationsgrößen. Es handelt sich um die R. Brauersche Gruppe der Algebrenklassen.

Für zyklische Algebren erhält man den Spezialfall: ist \(Z/k\) zyklisch und \(S\) eine erzeugende Substitution, so entsprechen den Potenzen von \(S\) die Potenzen von \(u\), und es kommt
\[
        A=Z+uZ+\cdots+u^{n-1}Z,                              \tag{1'}
\]
\[
        zu=uz^S,\qquad u^n=\alpha,                            \tag{2'--3'}
\]
wobei
\[
        \alpha\in k^*,\qquad
        \bar\alpha=\alpha\,\Norm(c)\quad\text{für }v=uc .      \tag{4'--5'}
\]
Jedes Faktorensystem besteht hier also aus einem einzigen Element des Grundkörpers, Bezeichnung \(A=(\alpha,Z,S)\); die Einsklasse wird durch die Normen aus \(Z^*\) gegeben, und die Gruppe der Algebrenklassen ist isomorph zu \(k^*/\Norm(Z^*)\).

Die Formulierung des Hauptgeschlechtssatzes im Minimalen benutzt die Tatsache, daß die Relationen (2) bis (5) rein multiplikativ sind und somit zugleich die aus den Elementen der \(n\) Komplexe \(u_S K^*\) bestehende Erweiterungsgruppe \(\Gg^*\) von \(\Gg\) definieren:
\[
        \Gg^*=\{\,u_{S_1}K^*,\ldots,u_{S_n}K^*\,\}.            \tag{6}
\]
\(\Gg^*\) besitzt \(K^*\) als abelschen Normalteiler, während \(\Gg^*/K^*\simeq\Gg\). \(\Gg^*\) besteht aus der Gesamtheit der regulären Elemente \(g^*\), die \(K\) in sich transformieren, also \(g^{*-1}Kg^*=K\). Die Ringautomorphismen von \(K\) werden also durch alle und nur die Elemente von \(\Gg^*\) erzeugt. Denn erzeugt \(v=us\) einen inneren Automorphismus von \(K\), so liegt \(w\) mit allen Elementen von \(K\) vertauschbar, also in \(K^*\), da \(K\) maximaler kommutativer Unterring ist.

\paragraph{Hauptgeschlechtssatz im Minimalen.}
Erste Fassung: Jeder Gruppenautomorphismus von \(\Gg^*\), der Fortsetzung des identischen Automorphismus von \(K^*\) ist, ist ein innerer und wird durch ein Element aus \(K^*\) erzeugt. Mit anderen Worten: haben die Transformationsgrößen
\[
        \frac{c_S^{\,T}c_T}{c_{ST}}
\]
sämtlich den Wert Eins, so sind die \(c_S\) symbolische \((1-S)\)-te Potenzen,
\[
        c_S=b^{1-S}=\frac{b}{b^S}\qquad(S\in\Gg).
\]
Dritte Fassung: Die Gruppe \(\Gg\) besitzt nur eine einzige zu Faktorensystem Eins gehörige verschränkte Darstellungsklasse ersten Grades in \(K^*\).

Eine Darstellung \(u_S\mapsto C_S\) heißt verschränkt zu Faktorensystem \(a_{S,T}\), wenn
\[
        C_S^{\,T}C_T=C_{ST}a_{S,T}
\]
gilt. Zwei Darstellungen gehören zur selben Klasse, wenn
\[
        C_S=B^{-S}D_SB .
\]

Ich beweise die erste Fassung und zeige dann ihre Übereinstimmung mit der zweiten und dritten. Der Beweis beruht darauf, daß jeder Gruppenautomorphismus der angegebenen Art sich zu einem Ringautomorphismus von \(A\) fortsetzen läßt, der bekanntlich innerer Automorphismus ist. Seien bei diesem Automorphismus \(v_S\) die Bilder der \(u_S\). Dann bildet \(\sum v_SK\) wieder das verschränkte Produkt von \(\Gg\) mit \(K\) zum selben Faktorensystem; denn nach Voraussetzung bleiben alle Relationen (2) bis (4) erhalten. Die Zuordnung
\[
        \sum u_Sb_S\longmapsto \sum v_Sb_S
\]
stellt einen Ringautomorphismus dar, bei dem die Elemente von \(K\) fest bleiben. Er wird daher durch ein Element aus \(K^*\) erzeugt.

Der Übergang zur zweiten Fassung beruht darauf, daß wegen (2) die \(v_S\) von der Form \(v_S=u_Sc_S\) sein müssen; da auch die Faktorensysteme erhalten bleiben, müssen nach (5) die zugehörigen Transformationsgrößen gleich Eins werden. Umgekehrt zeigt (2) bis (5), daß die Substitution \(v_S=u_Sc_S\) einen Automorphismus der angegebenen Art erzeugt. Nach der ersten Fassung erhält man also
\[
        v_S=b^{-1}u_Sb=u_S\,\frac{b}{b^S},
        \qquad c_S=b^{1-S}.
\]
Im zyklischen Fall bedeutet die Voraussetzung insbesondere \(\Norm(c)=1\), und \(c=b^{1-S}\) ist der bekannte Satz. Die dritte Fassung ist unmittelbar gleichbedeutend mit der zweiten, denn die Voraussetzung sagt aus, daß \(u_S\mapsto c_S\) eine verschränkte Darstellung ersten Grades zu Faktorensystem Eins bildet; \(c_S=b^{1-S}\) bedeutet die Äquivalenz zur Einsdarstellung.

\subsection*{§2. Der Hauptgeschlechtssatz}

\paragraph{Vorbemerkungen über Klasseneinteilung.}
Im auf zyklische Relativkörper bezüglichen Hauptgeschlechtssatz der Klassenkörpertheorie treten bekanntlich zwei verschiedene Klasseneinteilungen auf: absolute Idealklassen im Oberkörper, Strahlklassen im Unterkörper. Dem entspricht im allgemeinen Fall, daß eine Klasseneinteilung der Ideale des Oberkörpers eine feinere Idealklasseneinteilung für die Faktorensysteme induziert.

Bedeutet zunächst \(\Ii\) die Gruppe aller Ideale aus \(K\), so ist, da \(\Gg\) Automorphismen in \(\Ii\) erzeugt, die Erweiterung mit \(\Gg\) durch die Relationen (2) bis (5) definiert und besteht aus den \(n\) Komplexen \(\{\,\cdots u_S\Ii\cdots\,\}\). Geht man in \(\Ii\) durch Äquivalenzfestsetzung zur Gruppe der absoluten Idealklassen über, so sind auch jetzt die \(n\) Komplexe \(u_S\bar\Ii\) definiert. Denn \(\Gg\) erzeugt auch Automorphismen der Klassengruppe, da die Hauptklasse in sich transformiert wird. Damit überträgt sich die von \(\Ii\) zu \(\bar\Ii\) führende Äquivalenzbeziehung von \(\{\,\cdots u_S\Ii\cdots\,\}\) auf \(\{\,\cdots u_S\bar\Ii\cdots\,\}\). Insbesondere werden die \(u_S\) äquivalent zu allen Produkten \(u_S(c_S)\), wo \((c_S)\) die Hauptideale durchläuft.

Fordert man nun, daß die Produktrelationen (3) eindeutig sind im Sinne dieser Äquivalenz, so heißt das, daß alle Transformationsgrößen aus Hauptidealen eindeutig äquivalent werden. Mit anderen Worten:

In dem System der \(n\) Komplexe \(\{\,\cdots u_S\bar\Ii\cdots\,\}\), wo \(\bar\Ii\) die Gruppe der absoluten Idealklassen aus \(K\) und \(H\) insbesondere die Hauptklasse bedeutet, sind die Produktrelationen (3) für die \(u_SH\) dann und nur dann eindeutig definiert, wenn in der induzierten Idealklasseneinteilung der Faktorensysteme die Hauptklasse alle Transformationsgrößen aus Hauptidealen enthält.

\paragraph{Definition.}
Die induzierte Idealklasseneinteilung der Faktorensysteme wird dadurch festgelegt, daß die Hauptklasse aus allen Hauptidealen \((a_{S,T})\) besteht, für die mit mindestens einem Basiselement \(a_{S,T}\) die Algebra
\[
        (a_{S,T},K,\Gg)
\]
an allen Verzweigungsstellen von \(K\) zerfällt.

Diese Hauptideale bilden bezüglich der Verknüpfung der Faktorensysteme eine Gruppe; sie umfaßt die Transformationsgrößen
\[
        \frac{(c_S)^{\,T}(c_T)}{(c_{ST})},
\]
denn die Faktorensysteme \(c_S^{\,T}c_T/c_{ST}\) erzeugen überall zerfallende Algebren.

\paragraph{Formulierung des Hauptgeschlechtssatzes.}
Der Satz wird, entsprechend dem Satz im Minimalen, in drei gleichbedeutenden Fassungen ausgesprochen.

Erste Fassung: Entsteht, bei Zugrundelegung der induzierten Idealklasseneinteilung der Faktorensysteme, durch die Substitution
\[
        v_S=u_S c_S
\]
ein Automorphismus der \(n\) Komplexe \(\{\,\cdots u_S\bar\Ii\cdots\,\}\), d. h. bestehen für \(v_S\) im Sinne dieser Klasseneinteilung dieselben Relationen (3), so ist dieser Automorphismus ein innerer und wird durch eine Idealklasse \(b\) erzeugt.

Zweite Fassung: Gehören die aus den Idealklassen \(c_S\) gebildeten Transformationsgrößen
\[
        \frac{c_S^{\,T}c_T}{c_{ST}}
\]
sämtlich der Hauptklasse der induzierten Klasseneinteilung der Faktorensysteme an, so bilden die Vektoren \(\{\,\cdots c_S\cdots\,\}\) aus Idealklassen das Hauptgeschlecht und die Klassen \(c_S\) sind symbolische \((1-S)\)-te Potenzen; d. h. es gibt eine Idealklasse \(b\) mit
\[
        c_S=b^{1-S}=\frac{b}{b^S}\qquad(S\in\Gg).
\]

Dritte Fassung: Bei Zugrundelegung der induzierten Idealklasseneinteilung für die Faktorensysteme besitzt die Gruppe \(\Gg\) nur eine einzige, zur Einsklasse der Faktorensysteme gehörige verschränkte Darstellungsklasse ersten Grades in der Idealklassengruppe.

Daß diese Fassungen gleichbedeutend sind, folgt wie in §1. Der Übergang von der ersten zur zweiten Fassung wird hier sogar einfacher, da der Automorphismus schon durch \(v_S=u_Sc_S\) erzeugt vorausgesetzt wird. Zur dritten Fassung ist nur zu bemerken, daß Darstellungen in der Idealklassengruppe nur multiplikativ definiert sind; die Aussage über Darstellungen ersten Grades erfährt dadurch keine Einschränkung.

\paragraph{Beweis des Hauptgeschlechtssatzes.}
Der Beweis der zweiten Fassung beruht auf zwei Hilfssätzen.

\paragraph{Hilfssatz 1.}
Ergeben die aus den Idealen \(c_S\) gebildeten Transformationsgrößen
\[
        \frac{c_S^{\,T}c_T}{c_{ST}}
\]
das Einheitsideal, so sind die \(c_S\) symbolische \((1-S)\)-te Potenzen:
\[
        c_S=b^{1-S}\qquad(S\in\Gg).
\]
Gleichbedeutend ist dies mit der ersten Fassung, wenn man dort den Automorphismus durch \(v_S=u_Sc_S\) erzeugt voraussetzt. Beweis: \(b=\sum c_S\) leistet das Verlangte, wobei die Summe die Modulsumme, also den größten gemeinsamen Teiler, bedeutet. Denn
\[
        b=\sum_S c_S,\qquad
        b^T=\sum_S c_S^{\,T},\qquad
        b^T c_T=\sum_S c_S^{\,T}c_T=\sum_S c_{ST}=b,
\]
also \(c_T=b^{1-T}\).

\paragraph{Hilfssatz 2.}
Zerfällt die Algebra
\[
        A=(a_{S,T},K,\Gg)
\]
an allen endlichen und unendlichen Verzweigungsstellen von \(K\), und bilden die Hauptideale \((a_{S,T})\) Transformationsgrößen von Idealen,
\[
        (a_{S,T})=\frac{c_S^{\,T}c_T}{c_{ST}},
\]
so zerfällt \(A\) überall, also schlechthin. Dann werden die \((a_{S,T})\) Transformationsgrößen von Hauptidealen,
\[
        (a_{S,T})=\frac{(d_S)^{\,T}(d_T)}{(d_{ST})}.
\]
Umgekehrt genügt jede überall zerfallende Algebra dieser Bedingung.

Der Beweis beruht auf dem bekannten Verhalten von \(A\) bei \(p\)-adischer Erweiterung. Es bedeute \(k_p\) die \(p\)-adische Erweiterung von \(k\), entsprechend \(K_\PP\) die \(\PP\)-adische von \(K\), wobei \(\PP\) ein in \(p\) aufgehendes Primideal bedeutet; \(A_p\) sei die durch Übergang von \(k\) zu \(k_p\) entstandene Erweiterung von \(A\). Weiter sei \(\mathfrak Z\) die Zerlegungsgruppe zu \(\PP\), und \(a_{\mathfrak Z}\) der zugehörige Ausschnitt des Faktorensystems. Dann gilt
\[
        A_p\sim (a_{\mathfrak Z},K_\PP/k_p,\mathfrak Z).
\]
Ist \(p\) insbesondere unverzweigt, also \(\mathfrak Z\) zyklisch, so handelt es sich um die ausgezeichnete Darstellung von \(A_p\) vermöge des unverzweigten Trägheitskörpers \(K_\PP/k_p\).

Zu zeigen ist, daß unter der Voraussetzung von Hilfssatz 2 \(A_p\) auch in diesem unverzweigten Fall zerfällt. Für unendliche Stellen wird hier \(K_\PP=k_p\), also zerfällt \(A_p\). Für endliches \(p\) folgt zunächst, daß das Faktorensystem \(a_{\mathfrak Z}\) einem aus Einheiten bestehenden Faktorensystem assoziiert ist. Beim Übergang zur normierten zyklischen Darstellung wird in \(A_p=(\alpha,K_\PP/k_p,R)\), mit \(R\) Erzeugender von \(\mathfrak Z\), das Faktorensystem durch die Relationen
\[
        u_{\PP_i}u_{\PP_j}=u_{\PP_{i+j}}\varepsilon_{\PP_i,\PP_j}
\]
gegeben. Setzt man \(u_R=u\), so erhält man
\[
        u^i=u_R^i\,\varepsilon_R\varepsilon_{R^2}\cdots \varepsilon_{R^{i-1}},
\]
also insbesondere \(u_E=\varepsilon_E\), wenn \(E\) die Ordnung von \(\mathfrak Z\) ist. Weil bei unverzweigtem \(K_\PP/k_p\) alle Einheiten Normen sind, folgt, daß \(A_p\) zerfällt. Also zerfällt \(A\) überall, mithin nach dem Satz über zerfallende Algebren schlechthin.

Mit anderen Worten: aus der Voraussetzung folgt, daß \(a_{S,T}\) Transformationsgrößen aus Elementen sind,
\[
        a_{S,T}=\frac{d_S^{\,T}d_T}{d_{ST}},
        \qquad d_S\in K^*.
\]
Abgeschwächt heißt dies: die \((a_{S,T})\) werden Transformationsgrößen von Hauptidealen. Ist umgekehrt \(A\) überall zerfallend, so werden auch die Hauptideale \((a_{S,T})\) Transformationsgrößen von Idealen an jeder Stelle, und \(A\) zerfällt tatsächlich um eine andere Fassung der Voraussetzung des Satzes über die zerfallenden Algebren.

Aus den beiden Hilfssätzen folgt unmittelbar der Beweis des Hauptgeschlechtssatzes. Bedeutet \(c_S\) je ein Ideal aus der Klasse \(c_S\), so sagt die Voraussetzung, nach Definition der induzierten Klasseneinteilung, daß die Transformationsgrößen der \(c_S\) Hauptideale \((a_{S,T})\) werden, die den Voraussetzungen von Hilfssatz 2 genügen. Also gibt es Hauptideale \((d_S)\), so daß
\[
        \frac{c_S^{\,T}c_T}{c_{ST}}
\]
durch \((d_S)\) zur Einheit wird. Die Ideale \(c_S/(d_S)\), die in derselben Klasse \(c_S\) liegen, erfüllen somit die Voraussetzung von Hilfssatz 1. Also existiert eine Idealklasse \(b\) mit
\[
        c_S=b^{1-S}.
\]
Das bedeutet, daß die Klassen \(c_S\) symbolische \((1-S)\)-te Potenzen der Klasse \(b\) sind.

Im zyklischen Fall geht der Hauptgeschlechtssatz in den bekannten Satz über, sobald die Hauptklasse der induzierten Klasseneinteilung bei Zugrundelegung der Normierung in die Gruppe derjenigen Hauptideale \((\alpha)\) übergeht, für die \(\alpha\) an allen Verzweigungsstellen Normenrest wird.

\begin{center}
(Eingegangen am 27. 10. 1932.)
\end{center}

\section*{42. Zerfallende verschränkte Produkte und ihre Maximalordnungen}

\begin{center}
\emph{Actualités scientifiques et industrielles 148 (1934), S. 5--15}
\end{center}

In dem einfachsten Fall der zerfallenden verschränkten Produkte,\footnote{Die Theorie der verschränkten Produkte ist im Anschluß an eine Vorlesung von mir wiedergegeben in Abschn. II von H. Hasse, \emph{Theory of cyclic algebras over an algebraic number field}, Trans. Amer. Math. Soc. 34 (1932). Ich benutze hier nur die ersten Grundbegriffe. Die Brandtschen Grundlagen sind bei Hasse und in dem folgenden Rahmen neu begründet worden.} d. h. derjenigen, die zerfallende Algebren ergeben und also zu Eins assoziierte Faktorensysteme besitzen, lassen sich die Verhältnisse weitgehend explizit verfolgen. Das gilt auch noch bei nichtgaloisschen Zerfällungskörpern, d. h. maximalen kommutativen Teilkörpern. Die hierher gehörigen einfachen algebraischen Tatsachen werden in §1 entwickelt. In §2 gebe ich, bei algebraischem Zahlkörper als Grundkörper, explizite Darstellungen aller Maximalordnungen und ihrer Ideale vermöge Moduln und Komplementärmoduln des zugrunde liegenden Zerfällungskörpers \(k\).\footnote{Diese Darstellungen waren mir lange bekannt; den Anstoß zu den anschließenden Überlegungen gab die Mitteilung eines Satzes von H. Hasse.}

Ich gebe weiter eine Einteilung in Gebiete, wobei in ein Gebiet alle solchen Maximalordnungen gerechnet werden, die gleichen Durchschnitt mit \(k\) haben. Dadurch entsprechen die Gebiete eineindeutig den Ordnungen aus \(k\); der Hauptordnung ist das Hauptgebiet zugeordnet. Die Maximalordnungen eines Gebiets werden ineinander transformiert durch Ideale, die aus den Moduln der zugehörigen Ordnung gebildet sind; insbesondere handelt es sich beim Hauptgebiet um die Erweiterungen der Ideale aus \(k\). In §3 gehe ich zu beliebigen einfachen Algebren über. Durch Verfeinerung der Einteilung in Gebiete, indem neben dem Durchschnitt auch die \(p\)-adischen Komponenten der Maximalordnungen an den Verzweigungsstellen übereinstimmen müssen, kann man die im zerfallenden Fall gewonnenen Sätze übertragen. Insbesondere gehen auch dann die Maximalordnungen eines Hauptgebiets durch Transformation mit Erweiterungen der Ideale aus \(k\) ineinander über.

Die Sätze über das Hauptgebiet finden sich bei Chevalley und Hasse;\footnote{C. Chevalley, Sur certains idéaux d'une algèbre simple, Abh. Math. Sem. Hamburg, 1934; H. Hasse, Über gewisse Ideale in einer einfachen Algebra, dieser Arbeit vorangehend.} Chevalley zeigt auch, daß ohne die auf die Verzweigungsstellen bezügliche verfeinerte Einteilung die Sätze im allgemeinen nicht mehr gelten.

\subsection*{§I. Matrizeneinheiten der zerfallenden verschränkten Produkte und Komplementärbasen der Zerfällungskörper}

In diesem Paragraphen bedeute \(\Omega\) einen beliebigen Grundkörper. Es sei \(k/\Omega\) eine galoissche separable Erweiterung \(n\)-ten Grades, \(\Gg\) die Galoissche Gruppe mit den Elementen \(S,T,\ldots\); \(u_S,u_T,\ldots\) seien die zugeordneten Operatoren. Betrachtet werden verschränkte Produkte mit Faktorensystem Eins, also zerfallende Algebren
\[
        K=\Gg\times k=u_{S_1}k+\cdots+u_{S_n}k,\qquad
        u_Su_T=u_{ST},
\]
mit
\[
        z u_S=u_S z^S \qquad(z\in k).
\]
Setzt man
\[
        E_\Gg=\sum_{S\in\Gg}u_S,
\]
so erhält man aus dem Faktorensystem Eins die Relationen
\[
        E_\Gg u_S=u_S E_\Gg=E_\Gg,                            \tag{1}
\]
\[
        zE_\Gg=\sum_Su_Sz^S,                                  \tag{2a}
\]
und
\[
        E_\Gg z E_\Gg
        =E_\Gg\sum_Su_Sz^S
        =E_\Gg\,\Sp(z),                                      \tag{2b}
\]
wobei \(\Sp(z)\) die Spur von \(z\) als Element von \(k/\Omega\) bedeutet.

\paragraph{Satz.}
Bedeutet \(a_1,\ldots,a_n\) eine Basis von \(k/\Omega\) und \(\bar a_1,\ldots,\bar a_n\) die dazu komplementäre Basis, so ist durch
\[
        c_{ik}=\bar a_iE_\Gg a_k
\]
ein System von Matrizeneinheiten von \(K\) gegeben. \(K\) läßt sich als Modulprodukt darstellen in der Form
\[
        K=kE_\Gg k
          =k\Bigl(\sum_Su_S\Bigr)k
          =\sum_{i,k}(\bar a_iE_\Gg a_k)\,\Omega .
\]
Denn nach (2b) kommt
\[
        c_{ij}c_{hk}
        =\bar a_iE_\Gg a_j\bar a_hE_\Gg a_k
        =\bar a_iE_\Gg a_k\,\Sp(a_j\bar a_h)
        =c_{ik}\,\Sp(a_j\bar a_h).
\]
Dies ist gerade die Eigenschaft der Matrizeneinheiten, wegen der Definitionsgleichungen
\[
        \Sp(a_j\bar a_h)=\delta_{jh}
\]
der komplementären Basen.

\paragraph{Bemerkung 1.}
Ist \(Z\) irgendein kommutativer Erweiterungskörper von \(\Omega\), so bleibt alles erhalten, wenn \(a_i\) eine Basis von \(k_Z/Z\) bedeutet und \(K_Z\) direkte Summe von Körpern wird. Die Substitutionen der Gruppe sind dann dadurch definiert, daß sie auf \(Z\) die Identität induzieren.

\paragraph{Bemerkung 2.}
Daß \(K=kE_\Gg k\) ist, folgt auch unmittelbar daraus, daß die rechte Seite ein von Null verschiedenes Ideal des einfachen Systems \(K\) ist.

Sei nun \(\Omega\) von Charakteristik Null. Die Darstellung \(K=kE_\Gg k\) der zerfallenden Algebra bleibt gültig, auch wenn \(k\) ein nichtgaloisscher Zerfällungskörper von \(K/\Omega\) ist. Man geht zum galoisschen Zerfällungskörper \(\bar k\) über. Sei \(\bar k\) der zu \(k\) gehörige galoissche Körper, \(\bar\Gg\) seine Gruppe, \(\bar u_S\) die zugehörigen Operatoren, \(\bar K=\bar kE_{\bar\Gg}\bar k\) ein zerfallendes verschränktes Produkt. Der Unterkörper \(k\) gehöre zur Untergruppe \(\Hh\) der Ordnung \(h\). Man setzt
\[
        \bar\Gg=\Hh S_1+\cdots+\Hh S_n
              =S_1^{-1}\Hh+\cdots+S_n^{-1}\Hh,
        \qquad
        E_\Hh=\frac1h\sum_{H\in\Hh}u_H,
\]
und
\[
        E_\Gg=\frac1h E_{\bar\Gg}
              =\sum_i u_{S_i},\qquad
        u_{S_i}=E_\Hh\bar u_{S_i}.
\]
Dann erzeugt insbesondere das Idempotent \(E_\Hh\) die Einsdarstellung von \(\Hh\), und \(E_\Gg\) ebenso wie \(E_{\bar\Gg}\) die Einsdarstellung von \(\bar\Gg\). Für die so definierten \(E_\Gg\) und \(u_{S_i}\) gelten die Relationen (1) nach einer Seite und (2) vollständig. Es gilt nämlich zunächst
\[
        E_\Gg=E_\Gg E_\Hh=E_\Hh E_\Gg=E_\Hh E_\Gg E_\Hh,       \tag{3}
\]
\[
        zE_\Hh=E_\Hh z,                                      \tag{4}
\]
und folglich
\[
        z u_{S_i}=u_{S_i}z^{S_i}.                            \tag{5}
\]
Aus (5) folgt (2a) durch Summation, daraus (2b), und (1) folgt nach einer Seite:
\[
        E_\Gg u_{S_i}=E_\Gg .
\]
Der Beweis von \(K=kE_\Gg k\) erfolgt nun genau wie oben. Als voller Matrizenring über \(\Omega\) enthält \(K\) jeden Körper \(n\)-ten Grades als Unterkörper, also insbesondere auch \(k\).

Im Hinblick auf §3 sei noch bemerkt, daß der Übergang von \(K\) zu \(\bar K\) auch möglich ist, wenn \(K\) nicht zerfällt. Neben \(k\) ist dann auch \(\bar k\) Zerfällungskörper, also ist \(K\) ähnlich \(\bar K\), dem verschränkten Produkt aus \(\bar\Gg\) mit \(\bar k\). Dabei entspricht, da \(k\) Zerfällungskörper ist, der zu \(k\) gehörigen Untergruppe \(\Hh\) das Faktorensystem Eins; das Idempotent \(E_\Hh\) ist also wie oben definiert. Daraus ergibt sich rückwärts, daß \(K\) isomorph zu \(E_\Hh\bar K E_\Hh\) ist. Diese Algebra ist dem Automorphismenring des Linksideals \(\bar K E_\Hh\) isomorph, also zu \(\bar K\) ähnlich; der Rang stimmt mit dem von \(K\) überein. Für den zerfallenden Fall erhält man
\[
        E_\Hh\bar K E_\Hh
        =E_\Hh\bar kE_{\bar\Gg}\bar kE_\Hh
        =(E_\Hh\bar kE_\Hh)E_{\bar\Gg}(E_\Hh\bar kE_\Hh)
        =kE_\Gg k .
\]

\subsection*{§II. Die Maximalordnungen der zerfallenden verschränkten Produkte und ihre Einteilung in Gebiete}

Von jetzt an sei \(\Omega\) ein endlicher algebraischer Zahlkörper, \(K\) eine zerfallende Algebra \(n\)-ten Grades über \(\Omega\), und \(k\) ein galoisscher oder nichtgaloisscher maximaler kommutativer Teilkörper, also \(K=kE_\Gg k\). Es bedeute \(a,\bar a\) ein Paar komplementärer endlicher \(\oo\)-Moduln aus \(k\), wobei \(\oo\) die Hauptordnung von \(\Omega\) und \(o\) die Hauptordnung von \(k\) ist. Die Struktur der Maximalordnungen aus \(K\) relativ zu \(k\) ist durch die folgenden Sätze beschrieben.

\paragraph{Satz 1.}
Das Modulprodukt
\[
        \mathcal O_a=\bar a E_\Gg a
\]
ist eine Maximalordnung aus \(K\); insbesondere wird
\[
        \mathcal O=\mathcal O_o=oE_\Gg o
\]
Maximalordnung.

\paragraph{Satz 2.}
Die Transformation von \(\mathcal O\) in \(\mathcal O_a\) geschieht mittels des Linksideals
\[
        \mathcal L=oE_\Gg a
\]
und des dazu reziproken Rechtsideals
\[
        \mathcal R=\bar a E_\Gg o
\]
nach \(\mathcal O\).

\paragraph{Bemerkung.}
Das Zusammentreffen der beiden komplementären Moduln \(a,\bar a\) im Produkt \(\mathcal L\mathcal R\) führt gerade wegen der durch \(E_\Gg\) bewirkten Spurbildung zur Reziprokenbeziehung.

\paragraph{Satz 3.}
Alle Links- und Rechtsideale nach \(\mathcal O\) werden durch die in Satz 2 angegebenen erschöpft. Die Maximalordnungen \(\mathcal O_a\) erschöpfen also die Gesamtheit aller Maximalordnungen aus \(K\). Irgend zwei Maximalordnungen \(\mathcal O_a\) und \(\mathcal O_b\) gehen durch Transformation mit
\[
        \mathcal L_{ab}=\bar aE_\Gg b,
        \qquad
        \mathcal R_{ba}=\bar bE_\Gg a
\]
ineinander über; diese \(\mathcal L_{ab}\) und \(\mathcal R_{ba}\) erschöpfen die Links- und Rechtsideale nach \(\mathcal O_a\).

\paragraph{Satz 4.}
Durchläuft \(\mathcal O_a\) alle Maximalordnungen aus \(K\), so durchläuft der Durchschnitt
\[
        [\mathcal O_a,k]
\]
alle Ordnungen aus \(k\); dieser Durchschnitt ist jeweils gleich der Ordnung von \(a\). Faßt man alle zur selben Ordnung in \(k\) gehörenden Maximalordnungen zu einem Gebiet zusammen, so geschieht die Transformation innerhalb eines Gebiets durch Linksideale \(\bar aE_\Gg b\) nach \(\mathcal O_a\) und reziproke Rechtsideale, wobei \(a\) und \(b\) Moduln der zugehörigen Ordnung in \(k\) sind.

\paragraph{Satz 4a.}
Insbesondere geschieht die Transformation der Maximalordnungen des Hauptgebiets, also des Gebiets aller die Hauptordnung \(o\) von \(k\) enthaltenden Maximalordnungen, durch Ideale
\[
        aE_\Gg b,
\]
wo \(a\) und \(b\) Ideale aus \(k\) sind. Mit anderen Worten: ist \(\mathcal L\) Linksideal einer Maximalordnung des Hauptgebiets und ist \(o\) in der Rechtsordnung von \(\mathcal L\) enthalten, so ist \(\mathcal L\) Erweiterung eines \(o\)-Ideals.

Die Beweise beruhen auf dem Übergang zu den einzelnen Stellen. Es genügt jeweils der Übergang zum Quotientenring nach \(p\), wobei \(p\) alle Primideale aus \(\Omega\) durchläuft. Ist \(\mathcal C\) ein beliebiger \(\oo\)-Modul aus \(K\), so werde als Komponente \(\mathcal C_p\) der Erweiterungsmodul \(\mathcal C\oo_p\) bezeichnet. Dann gilt:

\paragraph{Hilfssatz.}
\(\mathcal C\) ist der Durchschnitt aller Erweiterungsmoduln \(\mathcal C_p\). Denn liegt \(w\) in \(\mathcal C_p\), so gibt es ein durch \(p\) nicht teilbares \(\sigma\in\oo\), so daß \(\sigma w\in\mathcal C\). Liegt \(w\) im Durchschnitt aller \(\mathcal C_p\), und bezeichnet \(g\) das Ideal aller \(\sigma\in\oo\) mit \(\sigma w\in\mathcal C\), so kann \(g\) durch kein \(p\) teilbar sein; also \(g=\oo\). Der Hilfssatz setzt den Höchstrang des Moduls nicht voraus.

\paragraph{Folgerungen.}
Jeder Modul ist durch seine Komponenten bestimmt. Ist \(\mathcal D\) echter Obermodul von \(\mathcal C\), so ist mindestens ein \(\mathcal D_p\) echter Obermodul von \(\mathcal C_p\). Ist insbesondere \(\mathcal M\) eine Ordnung in \(K\) und sind alle Komponenten \(\mathcal M_p\) Maximalordnungen nach \(\oo_p\), so ist auch \(\mathcal M\) Maximalordnung.

\paragraph{Beweis von Satz 1.}
\(\mathcal O_a\) ist eine Ordnung, denn \(\mathcal O_a\) ist endlicher \(\oo\)-Modul von Höchstrang. Weiter ist
\[
        \mathcal O_a^2
        =\bar aE_\Gg a\,\bar aE_\Gg a
        =\bar aE_\Gg a\,\Sp(a\bar a)
        \subseteq \bar aE_\Gg a=\mathcal O_a.
\]
Um die Maximalität zu zeigen, genügt es nach der Folgerung, die Maximalität jeder Komponente zu zeigen. Da \(\oo_p\) Hauptidealring ist, besitzen \(a_p\) und \(\bar a_p\) unabhängige komplementäre Basen \(a_1,\ldots,a_n\) und \(\bar a_1,\ldots,\bar a_n\). Die Produkte \(\bar a_iE_\Gg a_k\) bilden nach §1 Matrizeneinheiten; also wird
\[
        \mathcal O_{a,p}=\sum_{i,k}(\bar a_iE_\Gg a_k)\oo_p
\]
maximal. Da \(\mathcal O_a\) maximal ist, folgt auch \(\Sp(a\bar a)=\oo\).

\paragraph{Beweis von Satz 2.}
Aus der Spurenregel folgt:
\[
\begin{aligned}
 \mathcal O\mathcal L&=oE_\Gg o\,oE_\Gg a=oE_\Gg a=\mathcal L,\\
 \mathcal R\mathcal O&=\bar aE_\Gg o\,oE_\Gg o=\bar aE_\Gg o=\mathcal R,\\
 \mathcal L\mathcal R&=oE_\Gg a\,\bar aE_\Gg o=oE_\Gg o=\mathcal O,\\
 \mathcal R\mathcal L&=\bar aE_\Gg o\,oE_\Gg a=\bar aE_\Gg a=\mathcal O_a.
\end{aligned}
\]

\paragraph{Beweis von Satz 3.}
Sei \(\mathcal L\) ein Linksideal nach \(\mathcal O\). Dann ist insbesondere
\[
        \mathcal L=oE_\Gg\mathcal L
        =(oE_\Gg l_1,\ldots,oE_\Gg l_n)
\]
mit \(l_1,\ldots,l_n\) einer \(\oo\)-Basis von \(\mathcal L\). Setzt man \(l_\mu=\sum_i u_{S_i}z_{i\mu}\), so erhält man \(E_\Gg l_\mu=E_\Gg a_\mu\), also \(\mathcal L=oE_\Gg a\). Bei regulärem \(\mathcal L\) muß \(a\) Höchstrang \(n\) haben; dann ist \(\bar aE_\Gg o\) das reziproke Rechtsideal. Durch Zusammensetzung folgt die allgemeine Behauptung.

\paragraph{Beweis von Satz 4.}
Der Durchschnitt \(t=[k,\mathcal O_a]\) ist eine Ordnung. Als Untermenge von \(\mathcal O_a\) wird \(t\) Rechtsmultiplikatorenbereich und ist die größte Ordnung aus \(o\), die dieser Bedingung genügt; sie enthält also sicher die Ordnung \(o_a\) von \(a\). Um \(t=o_a\) zu zeigen, genügt Komponentengleichheit. Ist \(\tau\in t_p\), so ist \(\tau c_{ik}\) für jedes \(c_{ik}=\bar a_iE_\Gg a_k\) eine Linearkombination der \(c_{ij}\) mit Koeffizienten aus \(\oo_p\). Andererseits wird \(a_k\tau\) Linearkombination der \(a_j\) mit Koeffizienten aus \(\oo_p\). Wegen der Eindeutigkeit der Darstellung durch die \(c_{ij}\) müssen diese Koeffizienten in \(\oo_p\) liegen; also liegt \(\tau\) in der Ordnung von \(a_p\). Damit ist \(t\) gleich der Ordnung von \(a\).

\paragraph{Beweis von Satz 4a.}
Der Satz ist die Spezialisierung von Satz 4 auf das Hauptgebiet. Ist zunächst \(\mathcal O=oE_\Gg o\) und \(o\) auch in der Rechtsordnung von \(\mathcal L=oE_\Gg a\) enthalten, so gehört \(\bar aE_\Gg a\) zum Hauptgebiet; also ist \(a\) ein Ideal aus \(k\), und \(\mathcal L=\mathcal O a\). Geht man von \(\mathcal O_a\) aus, so wird entsprechend \(\mathcal L=\bar aE_\Gg b\), und \(a^{-1}b\) ist ein Ideal aus \(k\); es kommt \(\mathcal L=\mathcal O_a a^{-1}b\).

\subsection*{§III. Maximalordnungen beliebiger verschränkter Produkte}

Die für zerfallende verschränkte Produkte abgeleiteten Sätze bleiben auch im allgemeinen Fall bestehen, sobald man die Einteilung in Gebiete verfeinert. Die Tatsachen über das Hauptgebiet gelten dann genau, die übrigen müssen als Aussagen über die einzelnen Stellen formuliert werden. Dabei ist die Stelle jetzt als \(p\)-adische Erweiterung zu verstehen.

\paragraph{Definition.}
Sei \(K\) eine einfache normale Algebra über \(\Omega\) und \(k\) ein maximaler kommutativer Teilkörper. In ein Gebiet nach \(k\) werden alle Maximalordnungen aus \(K\) gerechnet, die denselben Durchschnitt mit \(k\) haben und außerdem an den Verzweigungsstellen von \(K\) dieselben \(p\)-adischen Komponenten besitzen. Ist der Durchschnitt mit \(k\) die Hauptordnung, so handelt es sich um ein Hauptgebiet.

Für die Hauptgebiete gilt dann:

\paragraph{Satz 1.}
Die Maximalordnungen eines Hauptgebiets gehen ineinander über durch Transformation mit Erweiterungsidealen der Ideale aus \(k\). Mit anderen Worten: die zur Differente primen Linksideale der Maximalordnungen eines Hauptgebiets, die \(o\) in ihrer Rechtsordnung enthalten, sind Erweiterungen der Ideale aus \(k\).

\paragraph{Satz 2.}
Ist \(K\) von Primzahlgrad, so gibt es nur ein Hauptgebiet; dieses geht durch Erweiterungsideale der \(k\)-Ideale in sich über. Die entsprechenden Ideale einer festen Maximalordnung \(\mathcal O\) sind die mit den zweiseitigen Idealen aus \(\mathcal O\) erweiterten Ideale aus \(k\).

\paragraph{Beweis.}
Zwei Maximalordnungen eines Hauptgebiets sind nur an endlich vielen Stellen verschieden; nach Definition sind dies Stellen, an denen \(K\) zerfällt. Hier wird aber \(K\) von der in §§1--2 betrachteten Form. Das ist klar, wenn \(k\) galoissch ist; wenn das Faktorensystem zu Eins assoziiert ist, kann es durch das Einssystem ersetzt werden. Nach §1 gilt dies auch im nichtgaloisschen Fall. Es handelt sich also an jeder solchen Stelle und damit allgemein um Erweiterungsideale.

Ist \(K\) von Primzahlgrad, so ist es entweder zerfallend oder Divisionsalgebra. Im zweiten Fall bleibt es an allen Verzweigungsstellen Divisionsalgebra und besitzt also nur ein Hauptgebiet, so daß wiederum Transformation mit Erweiterungsidealen gilt. Die allgemeinste Transformation erhält man durch Zusammensetzen der Erweiterungsideale mit den Idealen, die die Maximalordnung in sich transformieren, also mit den zweiseitigen Idealen; diese unterscheiden sich bekanntlich nur an den Verzweigungsstellen von Zentrumsidealen, also von Erweiterungsidealen.

Ferner gilt:

\paragraph{Satz 3.}
Die Maximalordnungen eines beliebigen Gebiets gehen ineinander über durch zur Differente prime Ideale, die an jeder Stelle aus Moduln der zugehörigen Ordnung zusammengesetzt sind, wie in §2, Satz 3 angegeben. Die Komponenten dieser Moduln bei einem festen Ideal sind nur an endlich vielen Stellen von der Ordnung verschieden.

Dies folgt unmittelbar aus §2; es ist nur zu beachten, daß die Operatoren \(u_S\) an den einzelnen Stellen durch entsprechende äquivalente ersetzt werden müssen.

Ob es auch hier zu jeder Ordnung aus \(k\) Maximalordnungen mit dieser Ordnung als Durchschnitt gibt, bedarf genauerer Untersuchung an den Verzweigungsstellen. Ein Hauptgebiet existiert stets, wie aus der verschränkten, eventuell verallgemeinerten Produktdarstellung entnommen werden kann, die zu einer \(o\) enthaltenden Ordnung führt, sobald die Faktorensysteme als ganz angenommen werden.

Eine weitere Frage ist, inwieweit es möglich ist, durch die arithmetische Einteilung nach Gebieten die algebraischen Zerfällungskörper zu charakterisieren. Mit anderen Worten: erzeugen verschiedene Zerfällungskörper stets verschiedene Gebietseinteilungen, erzeugen sie schon verschiedene Hauptgebiete? Erschöpfen etwa die Hauptgebiete aller Zerfällungskörper die Gesamtheit aller Maximalordnungen?

Daß durch eine einzelne Maximalordnung eine solche Charakterisierung sicher nicht möglich ist, zeigen die einfachsten Beispiele. So enthält die Maximalordnung
\[
        \left[\,1,i,j,\frac{1+i+j+k}{2}\,\right]
\]
des Quaternionenkörpers neben der Hauptordnung \([1,i]\) auch noch die Hauptordnung
\[
        \left[\,1,\frac{1+i+j+k}{2}\,\right]
\]
des Körpers der dritten Einheitswurzeln.

\begin{center}
Göttingen, August 1932.
\end{center}

\end{document}
